%! AP Statistics — Unit 7, Lesson 7.1 — Lesson Plan
\documentclass[10pt]{article}
\usepackage{apstats-article}
\usepackage{apstats-boxes}

\newcommand{\UnitNumberName}{Unit 7: Inference for Quantitative Data: Means \quad}
\newcommand{\LessonNumberName}{Lesson 7.1: Introducing Statistics --- Why Should I Worry About Error?}

\begin{document}
\begin{center}
    {\Large\bfseries \CourseName: \SchoolYear \\
    {\small\bfseries \UnitNumberName \LessonNumberName}}
\end{center}

\begin{tcolorbox}[colback=sky, colframe=navy, boxrule=0.9pt, arc=2mm,
    left=2mm, right=2mm, top=1.5mm, bottom=1.5mm]
    \textbf{Primary Objective:} Students will recognize that random variation can lead to
    errors in statistical inference, and they will frame the central question of Unit 7:
    how do we account for estimation error when making decisions about population means?
    (Big Idea: VAR)
\end{tcolorbox}

\renewcommand{\arraystretch}{1.6}
\begin{skillbox}[Priority Ideas \& Skills]{goldbox}
\begin{minipage}[t]{0.48\linewidth}
    \textbf{AP Skill 1 --- Selecting Statistical Methods}
    \begin{itemize}
        \item Identify questions suggested by probabilities of errors in statistical
              inference (Skill 1.A; VAR-1.I).
        \item Bridge: Unit 6 built inference for proportions using the normal model
              --- Unit 7 extends to population \emph{means}, where $\sigma$ is unknown.
        \item VAR-1.I.1: Random variation may result in errors in statistical
              inference --- wrong decisions can occur even when procedures are correct.
    \end{itemize}
\end{minipage}
\hfill
\begin{minipage}[t]{0.48\linewidth}
    \textbf{Key Understandings}
    \begin{itemize}
        \item Even a perfectly conducted inference procedure can yield an incorrect
              conclusion due to sampling variability.
        \item Knowing how often errors occur helps evaluate the reliability of our
              conclusions.
        \item This lesson is a framing lesson (Topic 7.1 is not directly assessed
              in the CED) that sets the stage for Lessons 7.2--7.10.
        \item Key shift from Unit 6: for means we estimate $\sigma$ with $s$,
              which produces the $t$-distribution instead of $z$.
    \end{itemize}
\end{minipage}
\end{skillbox}

\begin{skillbox}[{Vocabulary, Concepts, \& Theorems}]{greenbox}
\begin{minipage}[t]{0.48\linewidth}
    \begin{itemize}
        \item \textbf{Error in inference} --- reaching an incorrect conclusion about a
              population parameter because of sampling variability
        \item \textbf{Type I error} --- rejecting a true null hypothesis (false positive);
              probability = $\alpha$
        \item \textbf{Type II error} --- failing to reject a false null hypothesis
              (false negative); probability = $\beta$
    \end{itemize}
\end{minipage}
\hfill
\begin{minipage}[t]{0.48\linewidth}
    \begin{itemize}
        \item \textbf{Significance level ($\alpha$)} --- the maximum tolerable probability
              of a Type I error
        \item \textbf{Population SD ($\sigma$)} --- the true spread of the population;
              rarely known in practice
        \item \textbf{Sample SD ($s$)} --- our estimate of $\sigma$, computed from data;
              replaces $\sigma$ in Unit 7 procedures
        \item \textbf{$t$-distribution} --- family of distributions arising when $s$
              replaces $\sigma$; heavier tails than the normal
    \end{itemize}
\end{minipage}
\end{skillbox}

\begin{skillbox}[Activate Prior Knowledge \& Spiral Review]{sky}
\begin{minipage}[t]{0.48\linewidth}
    \textbf{Prior Knowledge Check (3--4 min)}
    \begin{itemize}
        \item Unit 6: one-sample $z$-interval and $z$-test for a proportion.
        \item Unit 6: Type I and Type II errors with proportions --- what each means
              in a real scenario.
        \item Unit 5: sampling distribution of $\bar x$ --- shape (CLT), center
              $\mu_{\bar x} = \mu$, spread $\sigma_{\bar x} = \sigma/\sqrt{n}$.
        \item Connect: ``What happens when $\sigma$ is unknown?''
    \end{itemize}
\end{minipage}
\hfill
\begin{minipage}[t]{0.48\linewidth}
    \textbf{Connections Forward}
    \begin{itemize}
        \item $t$-distribution and its properties $\to$ Lesson 7.2
        \item Confidence intervals for one mean $\to$ Lessons 7.2--7.3
        \item Significance tests for one mean $\to$ Lessons 7.4--7.5
        \item Two-sample $t$-procedures $\to$ Lessons 7.6--7.9
        \item Choosing between $z$ and $t$ $\to$ Lesson 7.10
    \end{itemize}
\end{minipage}
\end{skillbox}

\begin{skillbox}[Hook \& Lesson]{sky}
\begin{minipage}[t]{0.48\linewidth}
    \textbf{Hook (5--7 min)}\\
    Display a headline: ``Drug lowers average recovery time by 2 days --- results
    statistically significant.'' A follow-up: ``New study finds no effect.''\\[4pt]
    Ask: ``How can two studies disagree? Did someone make a mistake?''\\[4pt]
    Reveal: Both studies may have followed correct procedures.
    \textbf{Random variation means even valid inference can lead to wrong conclusions.}
    The question of Unit 7: \textbf{How do we quantify and manage error when working
    with population means?}
\end{minipage}
\hfill
\begin{minipage}[t]{0.48\linewidth}
    \textbf{Lesson Flow (15 min)}
    \begin{enumerate}
        \item Revisit the sampling distribution of $\bar x$ from Unit 5. Shape:
              approximately normal by CLT when $n \ge 30$ (or population is normal).
              Center: $\mu$. Spread: $\sigma/\sqrt{n}$.
        \item The catch: $\sigma$ is almost never known. Substituting $s$ introduces
              additional variability --- that is why the $t$-distribution has heavier
              tails than the standard normal.
        \item Foreshadow: when we use $s$, the statistic
              $t = (\bar x - \mu)/(s/\sqrt{n})$ follows a $t$-distribution with
              $n - 1$ degrees of freedom.
        \item Frame errors in the means context: a Type I error means concluding
              $\mu \ne \mu_0$ when it is actually equal; a Type II error means failing
              to detect a real shift in the mean.
    \end{enumerate}
\end{minipage}
\end{skillbox}

\begin{skillbox}[Explicit Instruction \& Active Monitoring]{sky}
\begin{minipage}[t]{0.48\linewidth}
    \textbf{Direct Instruction (8 min)}
    \begin{itemize}
        \item Show a simulation: 100 samples from a population with known $\mu = 50$
              each produce a 95\% CI for $\mu$. Roughly 5 of the 100 intervals miss
              $\mu = 50$ --- that is a Type I error rate of 5\%.
        \item Connect $\alpha$ to error rate: if $\alpha = 0.05$, we expect to commit
              a Type I error in about 5\% of correctly conducted tests.
        \item Emphasize: error rates are properties of the \emph{procedure}, not of
              any individual sample. A single result is either right or wrong --- we
              just cannot know which.
    \end{itemize}
\end{minipage}
\hfill
\begin{minipage}[t]{0.48\linewidth}
    \textbf{Active Monitoring}
    \begin{itemize}
        \item Listen for ``statistical significance means no error.'' Redirect:
              significance only means the result is unlikely under $H_0$; Type I
              errors still occur at rate $\alpha$.
        \item Cold-call: ``If $\alpha = 0.01$ instead of 0.05, how does that change
              the Type I error rate? The Type II error rate?''
        \item Check warm-up for fluency with the sampling distribution of $\bar x$
              and with identifying error types from Unit 6 scenarios.
    \end{itemize}
\end{minipage}
\end{skillbox}

\begin{skillbox}[Group Work \& Differentiation]{redbox}
\begin{minipage}[t]{0.48\linewidth}
    \textbf{Partner Activity (8--10 min)}\\
    Three real-world decision scenarios (drug trial, quality control, employee survey).
    For each:
    \begin{enumerate}
        \item State the null and alternative hypothesis in words.
        \item Describe a Type I error in context.
        \item Describe a Type II error in context.
        \item Decide which error is more costly and explain why.
    \end{enumerate}
\end{minipage}
\hfill
\begin{minipage}[t]{0.48\linewidth}
    \textbf{Differentiation}
    \begin{itemize}
        \item \textit{Support (Tier R)}: Provide sentence frames for error descriptions;
              supply the hypotheses already stated.
        \item \textit{On-grade (Tier A)}: Complete all four steps for each scenario.
        \item \textit{Extension (Tier E)}: For one scenario, argue whether lowering
              $\alpha$ from 0.05 to 0.01 is a good trade-off. What happens to Type II
              error and power?
        \item \textit{ELL}: Bilingual card --- ``Type I error / error tipo I,''
              ``significance level / nivel de significancia.''
    \end{itemize}
\end{minipage}
\end{skillbox}

\begin{skillbox}[Individual Work \& Assessment]{redbox}
\begin{minipage}[t]{0.48\linewidth}
    \textbf{Exit Ticket (5 min)}
    \begin{enumerate}
        \item A quality-control inspector tests whether the mean fill of juice bottles
              differs from the labeled 12 oz. Describe in context what a Type I error
              would mean.
        \item Describe what a Type II error would mean in the same context.
        \item If the inspector sets $\alpha = 0.10$ instead of 0.05, how does that
              change the Type I and Type II error rates?
    \end{enumerate}
\end{minipage}
\hfill
\begin{minipage}[t]{0.48\linewidth}
    \textbf{AP-Style Check}\\
    \textit{A researcher sets $\alpha = 0.05$ and conducts a significance test. She
    fails to reject $H_0$. Which of the following is a correct interpretation?}
    \begin{enumerate}[label=(\Alph*)]
        \item She has proved that $H_0$ is true.
        \item The probability that $H_0$ is true is 0.95.
        \item She may have committed a Type II error.
        \item She has committed a Type I error.
    \end{enumerate}
    Answer: (C)
\end{minipage}
\end{skillbox}

\begin{skillbox}[Reinforcement \& Extension]{goldbox}
\begin{minipage}[t]{0.48\linewidth}
    \textbf{Homework / Practice}
    \begin{itemize}
        \item For three given scenarios (each with a context, $H_0$, and $H_a$):
              (a) describe a Type I error in context; (b) describe a Type II error
              in context; (c) state which error is more consequential and justify.
        \item Reflection: ``Explain why it is impossible to simultaneously minimize
              both the Type I and Type II error rates for a fixed sample size.''
    \end{itemize}
\end{minipage}
\hfill
\begin{minipage}[t]{0.48\linewidth}
    \textbf{Extension / Preview}
    \begin{itemize}
        \item \textit{Extension}: Research a published clinical trial that was later
              contradicted. Identify whether the initial or follow-up study was more
              likely to have committed a Type I or Type II error. Justify.
        \item \textit{Preview}: Lesson 7.2 introduces the $t$-distribution and builds
              the first complete inference procedure for a population mean --- the
              one-sample $t$-interval.
    \end{itemize}
\end{minipage}
\end{skillbox}

\end{document}
